% 正例样本：合规写法参考（检查器应零命中）

\section{问题背景}

某市 2024 年早高峰主干道平均车速为 18.3 km/h，较 2019 年的 24.1 km/h 下降 24.1\%（来源：市交通委年度报告）。
既有模型在样本量小于 500 时 MAPE 普遍高于 15\%，本题可用的历史记录仅 320 条，因此需要小样本策略。

\section{模型建立与求解}

本文以最小化预测误差为目标建立梯度提升模型，并用 5 折交叉验证选择超参数。求解分两步：先用滑动窗口
构造滞后特征与时段特征，再以 TPE 搜索树深、学习率与子采样比例。最终 MAPE 为 8.7\%，较基准的 13.2\%
下降 4.5 个百分点；在 ±10\% 的流量扰动下 MAPE 波动不超过 0.6 个百分点。

\section{结论}

本文给出的早高峰流量预测模型在 320 条样本上达到 MAPE 8.7\%，满足题面「误差不超过 10\%」的要求，
可直接用于信号配时方案的滚动评估；当样本量低于 150 条时该结论尚待验证。
